Tool-using large language model (LLM) agents convert natural-language requests into consequential external effects, so semantic uncertainty can become an authorization failure even when a tool call is syntactically valid. We present \system, a semantic-to-effect authorization framework that preserves a \emph{set} of plausible interpretations, compiles them into effect programs, and authorizes only their common safe core. A Semantic Anchor Lock independently preserves security-critical surface constraints, while a Clarification Evidence Gate recovers utility only after anchor-constrained interpretations are effect-equivalent to authenticated authorization.

On a frozen, group-disjoint internal holdout of 2,000 cases (500 semantic groups across English, Romanized Hindi, Romanized Telugu, and Telugu script), Qwen3.5-122B-A10B-FP8 reaches 99.85\% final decision accuracy after all 80 observed pre-guard unsafe \allow decisions are intercepted; GPT-OSS-120B reaches 99.80\% after both observed pre-guard unsafe \allow decisions are intercepted. Both finish with zero observed unsafe \allow decisions and zero false blocks. A 50,000-replicate group-cluster bootstrap gives 95\% intervals of 99.65--100.00\% and 99.60--99.95\%, respectively. Frozen-log ablation shows that Anchor Lock supplies the security gain, whereas clarification recovery supplies utility and is not safe as a standalone substitute for anchoring.

Post-freeze metamorphic testing exposed a specific \texttt{matching}$\rightarrow$\texttt{all} scope-broadening weakness in the frozen Anchor Lock v2. Anchor Lock v3 closes all 1,496 previously missed scope-broadening cases and achieves 100\% recall on 10,248 critical metamorphic attacks with zero false positives on 2,000 clean cases; it introduces no new triggers on 3,996 evaluable frozen model outputs. Separately, a transaction-first reference execution broker passes 1,320 distinct randomized replay/tamper/staleness/crash scenarios, but remains a reference prototype rather than an integrated production executor. A targeted hosted $\tau^3$ QUICK pilot contributes six benchmark-native trajectories across Qwen and GPT-OSS, with mean rewards 0.33 and 0.67. A subsequent local-T4 replication served multiple open-weight checkpoints and verified automatic tool calling on Granite-3.3-2B, Qwen2.5-3B, Qwen2.5-7B-AWQ, and xLAM-2-3B, but BFCL timed out without retained scores, AgentDojo completed only banking slices, and local $\tau^3$ remained partially evaluable because of context-window limits. We therefore retain a deliberately bounded claim: strong internal evidence with initial but incomplete external validation and production mediation still open.